\documentclass{jcgt}
\pdfminorversion=7

\setciteauthor{Malyshau}
\setcitetitle{Rendering Editable Layered Terrain with WebGPU: A Comparison of Six Methods}
\setheadtitle{Rendering Editable Layered Terrain with WebGPU}
\submitted{2026-08-17}
\lhead{\small Manuscript for the \textit{Journal of Computer Graphics Techniques}}
\rhead{\small Revised draft: 2026-10-05}

\usepackage{booktabs}
\usepackage{array}
\usepackage{tabularx}
\usepackage{placeins}
\usepackage{flafter}
\newcolumntype{Y}{>{\raggedright\arraybackslash}X}
\newcommand{\code}[1]{\texttt{#1}}
\newcommand{\methodfigure}[3]{%
  \begin{figure}[tbp]\centering
  \includegraphics[width=\linewidth]{figures/#1.pdf}
  \caption{#2}\label{#3}
  \end{figure}}
\captionsetup{font=small,skip=6pt}
\lstset{basicstyle=\small\ttfamily,columns=fullflexible,
  showstringspaces=false,breaklines=true,frame=single,language=}
\setlength{\emergencystretch}{1em}

\begin{document}

\title{Rendering Editable Layered Terrain with WebGPU:\\A Comparison of Six Methods}
\author{Dzmitry Malyshau~\href{https://orcid.org/0009-0005-6410-4276}{\includegraphics[width=8pt]{ORCIDlogo}}\\Independent Researcher}
\maketitle

\begin{abstract}
Terrain with caves and destructible ground needs more than a single height per map location. We compare six ways to render the layered terrain of \emph{Vangers}: height-field ray marching, voxel-assisted ray marching, horizontal slices, per-sample bars, point scattering, and a fitted triangle mesh. All six run in one engine through wgpu, using the same terrain data, camera, and shared shading code. We measure image quality, frame time, preparation, memory, and local edits on five devices.

Camera angle changes the tradeoffs. Point scattering leaves large gaps near the horizon, and horizontal slices produce bands. At the selected settings, the mesh has the lowest mean frame time on every measured device. It also retains 319 MiB of GPU geometry and 535 MiB of CPU triangulation. The direct ray marcher uses no extra persistent method data and reads edits on the next draw. A survey of ten worlds links poor mesh reduction to the second terrain layer; an insertion count shows that layer boundaries account for 23\% of vertex insertions in one Fostral fit. These results help choose a renderer for this terrain format and show why frame time alone is not enough. The engine and measurement tools are open source.
\end{abstract}

\section{Introduction}\label{sec:introduction}

A height field stores one surface height at each horizontal location. It is a useful model for terrain, but it cannot describe a cave beneath a solid roof. The terrain in \emph{Vangers} (1998) adds a second solid interval to each column. Players can also change the terrain during play. A renderer must therefore handle floors, roofs, and their boundary walls, then update them after an edit.

Several GPU methods can draw this data. Their tradeoffs are hard to judge from a single screenshot or timing result. A method that looks complete from above can leave holes at eye level. A mesh can draw quickly while taking time and memory to build. An acceleration structure can speed up a ray query while adding work after each edit.

We study these choices in vange-rs, an open-source engine for the Vangers terrain format. It contains six rendering methods built over a common data path. We ask three questions. How does the view direction affect image quality and frame time? What does each method retain and update? Why does a fitted mesh reduce some worlds much more than others?

The main result is a practical tradeoff between a fitted mesh and a direct ray marcher. The selected mesh is the fastest on average on all five tested devices, but needs a blocking initial fit, substantial retained data, and local refits. The ray marcher is slower at these settings, but adds no persistent acceleration data and reads terrain edits directly. Slices and points have more visible sampling failures near the horizon (Figure~\ref{fig:teaser}).

Our contribution is a comparison of six implementations under a shared workload, together with an analysis of their failure modes and update costs. We also study the cost of fitting a shared mesh to three altitude fields. The rendering families themselves are established techniques. We do not claim a new ray marcher, voxel hierarchy, splat method, or terrain triangulator.

Section~\ref{sec:data} defines the terrain and edit model. Section~\ref{sec:methods} explains the methods. Sections~\ref{sec:eval} and~\ref{sec:results} describe the study and its results. Section~\ref{sec:fit} examines mesh size, and Section~\ref{sec:choice} draws out the choice between methods.

\methodfigure{teaser}{The six selected configurations at the same eye-level hangar view. Scattered leaves gaps, and Sliced shows bands. The selected Mesh configuration preserves the wall. These images share the camera, palette, fog, and shadow pass.}{fig:teaser}

\section{Terrain and Edits}\label{sec:data}

The decoded terrain is a two-dimensional grid of vertical columns. A column contains ground up to an altitude \code{low}. It may also contain an upper slab, whose underside is \code{mid} and whose top is \code{high}. Thus a point in a dual-level column is solid when
\begin{equation}
  z \leq \code{low}(x,y)
  \quad\hbox{or}\quad
  \code{mid}(x,y) \leq z \leq \code{high}(x,y).
  \label{eq:solid}
\end{equation}
For a single-level column, \code{low = mid = high}, so the same test reduces to one solid interval. Figure~\ref{fig:encoding} shows the two cases.

The stored format packs a dual-level column into a pair of horizontal texels. The even texel stores the floor height, the odd texel stores the slab top, and metadata bits in both encode the ceiling offset. All methods use the same decoder. Heights have 8-bit precision. Fostral, the world used for the rendering comparison, has $2048\times16384$ texels; 10.9\% use the dual-level encoding.

\methodfigure{encoding}{Decoded terrain columns. A single-level column has one solid interval. A dual-level column adds a slab above the floor. The gap between \code{low} and \code{mid} is empty. Heights and layer presence can both change during play.}{fig:encoding}

This terrain is hand-authored. Adjacent columns can have large height jumps, and a slab can begin or end at a column boundary. Interpolating across such a boundary can create a sloping surface where the encoded terrain has a wall. This matters both when rendering the terrain and when fitting it with triangles.

A local edit changes heights or layer metadata in a dirty rectangle. The engine uploads the changed texture rows. Ray, Sliced, Painted, and Scattered then read the new values during their next draw. Voxel updates the affected occupancy cells and their ancestors. Mesh refits the affected chunks and replaces their GPU buffers. We measure these update costs separately from ordinary frames.

The workload requires all methods to represent both solid intervals, support local changes without a full level reload, and run within an interactive application. These are design requirements, not claims that every configuration meets every frame budget or image-quality target. We report the measured costs and failures.

\section{Rendering Methods}\label{sec:methods}

All methods share the decoded height and material data, camera, palette, fog, diffuse-lighting function, and shadow map. They differ in how they find visible terrain samples. Mesh also differs in its normal input: it can use the triangle normal on walls and ceilings, while the other methods use a height-field gradient. The comparison therefore includes this representation-dependent shading difference.

Table~\ref{tab:methods} groups the methods by the work they generate. Ray and Voxel start from screen pixels. Sliced, Painted, and Scattered generate terrain samples and project them to the screen. Mesh first builds a triangle surface, then draws it with ordinary rasterization. We use these short names throughout; the harness calls Ray \code{RayTraced} and Voxel \code{RayVoxel}.

\begin{table}[tbp]\centering\small
\caption{How each method draws the terrain and responds to an edit.}\label{tab:methods}
\begin{tabularx}{\linewidth}{@{}l Y Y@{}}
\toprule
Method & Work sent to the GPU & Work after a local edit \\
\midrule
Ray & One ray per pixel & Read changed texture \\
Voxel & One ray per pixel, occupancy hierarchy & Update occupancy hierarchy \\
Sliced & Horizontal planes & Read changed texture \\
Painted & Bars at ground samples & Read changed texture \\
Scattered & Points from a compute grid & Read changed texture \\
Mesh & Fitted triangles & Refit changed chunks \\
\bottomrule
\end{tabularx}
\end{table}

\subsection{Ray: sample the height field directly}\label{sec:ray}

A full-screen pass reconstructs a ray for each pixel. It clips the segment at the far plane or, for a downward ray, at $z=0$. It then advances with a fixed step and applies Equation~\ref{eq:solid}. The first solid sample brackets a crossing with the preceding sample. Four bisection steps refine that bracket. A miss leaves the cleared depth; a hit writes depth so that terrain composes with other geometry.

The selected setting uses 128 forward samples. A step can skip a feature thinner than the distance between samples, especially along a long, grazing ray. Bisection improves a detected crossing but cannot recover a feature that the forward pass skipped. Figure~\ref{fig:sampling} illustrates this distinction. The method has no acceleration structure to prepare or maintain. Appendix~\ref{sec:ray-code} gives pseudocode with the endpoint and miss handling used by the shader.

\subsection{Voxel: skip empty regions}\label{sec:voxel}

Voxel adds a conservative occupancy pyramid to the same terrain. The finest level stores occupancy bits in Morton-ordered $8\times8\times8$ tiles. Each coarser cell is the union of eight children. A hierarchical grid traversal skips an empty cell, descends into an occupied cell, and samples the original height field inside an occupied leaf. The hierarchy identifies empty space; it does not replace the surface with voxel faces.

The selected cell dimensions are $(4,8,2)$ world units. The outer traversal and inner sampling loops each have a budget of 100 steps, so dense or long rays can exhaust a budget. The hierarchy takes 18.29 MiB. The GPU builds it in batches and updates cells touched by an edit. This saves empty-space tests but adds memory and preparation work.

\subsection{Sliced: intersect horizontal planes}\label{sec:sliced}

Sliced draws horizontal quads over the camera's terrain footprint. A fragment survives only if its height lies inside a solid interval. Reducing the slice count increases the spacing over the same full altitude range. The selected setting uses 512 slices over 256 quantized height levels.

This method reads edits directly and needs no stored geometry. Its weakness is the gap between planes. At grazing angles, the gaps become visible as bands (Figure~\ref{fig:sampling}). Increasing the slice count reduces this error but also increases fragment work.

\methodfigure{sampling}{Sampling in a vertical cross-section. Top: a fixed-step ray first brackets the surface between an empty and a solid sample; bisection refines only that bracket. A thinner feature between forward samples could be missed. Bottom: slices retain fragments only inside the solids. The vertical space between slices can appear as bands from a grazing view.}{fig:sampling}

\subsection{Painted: draw a bar at each sample}\label{sec:painted}

Painted creates column geometry in the vertex shader. Each ground sample produces a bar from zero to \code{low}; a dual-level sample adds a bar from \code{mid} to \code{high}. The shader emits the three faces oriented toward the camera. Positions come from the instance index, so the CPU does not upload a vertex stream for each bar.

The draw covers an axis-aligned bound of the camera footprint. Samples are ordered from front to back so that early depth tests can reject covered fragments. Work still grows with the number of ground samples in the footprint. Edits need only the shared terrain upload.

\subsection{Scattered: project points with compute}\label{sec:scattered}

Scattered uses a compute grid over a camera-aligned ground footprint. The sampling grid is denser near the camera. Each invocation samples both solid intervals and projects the points into a screen buffer. A 32-bit \code{atomicMin} keeps the nearest packed value: 24 bits of depth and 8 bits of terrain type. A full-screen pass reconstructs the position and applies the shared shading code.

This approach follows the object-order projection of the original Vangers software renderer, with atomics replacing its painter ordering. Each sample writes one pixel. It has no filtered footprint, so projection does not guarantee coverage. Too few samples leave gaps; changes in which point wins can produce speckle. We evaluate this point scatterer, not a general surface-splatting system.

\subsection{Mesh: fit and retain triangles}\label{sec:mesh}

Mesh fits a triangulated irregular network (TIN) by greedy point insertion, following \citet{garland1995terrain}. Each triangle tracks the source sample with the largest vertical error. The fitter inserts that sample and restores the Delaunay triangulation. It repeats until the error falls below a tolerance. Integer grid coordinates allow exact orientation and incircle predicates for this data.

One planar topology serves the \code{low}, \code{mid}, and \code{high} fields. Any of the three can trigger an insertion. The renderer emits the floor, slab surfaces, and walls at dual-region edges. The fit is split into $128\times128$ chunks, with three levels of detail (LODs). The draw selects an LOD by distance and culls chunks outside the view. Shared chunk borders use a stable vertex sequence to avoid cracks.

The quality parameter $q$ sets the vertical tolerance to $16\,2^{-4q}$ height-quantization steps. One step is one world unit in this study, so $q=0.5$ gives a tolerance of 4 units. The renderer retains all three GPU LODs and the CPU triangulations used for local edits. A refit inserts vertices into the existing triangulation. Its result can therefore depend on earlier edits, unlike a fresh fit of the same final terrain.

\subsection{Relation to earlier methods}\label{sec:related}

Layered terrain models occupy the space between a single height field and a full volume. \citet{benes2001layered} use vertical material intervals for erosion. Grounded Heightmap Trees use locally oriented height maps to represent more general overhangs and tunnels \citep{alonso2008grounded}. QuadStack compresses layered data and ray casts it on the GPU \citep{graciano2021quadstack}. Here we keep the game's fixed two-interval format and require local gameplay edits. Layered Depth Images also store multiple depths, but along an input camera's rays rather than along world-space vertical columns \citep{shade1998ldi}.

The ray methods build on linear search and refinement in relief mapping, landscape ray casting, and conservative skipping \citep{policarpo2005relief,mantler2006landscape,tevs2008maximum}. Our occupancy pyramid uses a regular hierarchy rather than a compressed sparse voxel scene \citep{laine2010svo}. Sliced is related to plane-based volume rendering \citep{cabral1994volume}; Scattered is simpler than footprint-based volume or surface splatting \citep{westover1990footprint,zwicker2001splatting}.

TIN fitting and terrain LOD are well studied for single-valued height fields \citep{fowler1979tin,garland1995terrain,duchaineau1997roam,losasso2004clipmaps}. Our mesh uses one topology for three altitude fields, including sharp transitions where a slab ends. Section~\ref{sec:fit} measures the cost of that choice. This study compares the implementations described above; it does not rank the best possible implementation of each family.

\FloatBarrier
\section{Evaluation}\label{sec:eval}

The study has three parts: a twelve-view rendering comparison on Fostral, a local-edit experiment on the same world, and a ten-world survey of mesh fitting. This section defines the rendering and edit tests. Section~\ref{sec:fit} defines the fit survey.

\subsection{Scenes, devices, and timing}\label{sec:protocol}

The rendering test uses twelve fixed camera records: three views at each pitch of $0^\circ$, $-30^\circ$, $-60^\circ$, and $-90^\circ$. Zero pitch looks toward the horizon; $-90^\circ$ looks straight down. Each image is $1280\times800$ pixels, with a near distance of 1 and a far distance of 600 world units. Every method renders the same $1024\times1024$ height-field shadow map.

The five devices are an AMD Radeon 780M, AMD Radeon RX 7900 XT, Intel RPL-U GPU, NVIDIA GeForce RTX 5070, and Apple M3. All runs use native wgpu and the same WGSL shader sources \citep{wgpu,w3cwebgpu,w3cwgsl}. The first four use Vulkan; the M3 uses Metal. The measured code is pinned at commit \code{5586e2a}, tagged \code{terrain-paper}. Device and driver records are in \path{paper/results.md}.

After method-specific warmup, we time 40 frames per view. Each frame is submitted and awaited on its own. Vulkan times come from GPU timestamps. On Metal, the encoder-level timestamps did not reliably enclose the multipass workload, so we use CPU submit-and-wait times. The M3 results support comparisons between methods on that device; they are not directly comparable with the Vulkan GPU times. These isolated-frame measurements are latency measurements, not pipelined game frame rates.

We report the mean over the twelve views. Within-session 95\% intervals use the recorded frame samples, as described in Appendix~\ref{sec:timing-detail}. They do not measure variation between sessions or drivers. Preparation, memory, and edits are measured separately.

\subsection{Image-quality measures}\label{sec:metrics}

A CPU ray cast of the decoded terrain supplies a hit mask and a distance along each camera ray. It searches the segment with coarse samples and refines detected crossings by bisection. This is a sampled reference, not an exact ray--column intersection oracle. Its remaining limitations are stated below.

We use three measures:
\begin{itemize}
\item \textbf{Coverage.} Missing terrain is the fraction of all image pixels where the reference hits terrain but the renderer leaves the background. Added terrain is the fraction where the renderer hits but the reference sees background. The harness calls these \code{see\_through} and \code{covers\_sky}. Both use the full image as the denominator.
\item \textbf{Distance error.} Where both images hit terrain, we report the median and 95th percentile of the absolute difference in ray distance, in world units.
\item \textbf{Local incoherence.} We flag a hit pixel when its distance differs from the median of its $3\times3$ neighborhood by more than $0.04\max(d,1)$, where $d$ is that pixel's distance. Cleared pixels have undefined distance; if the neighborhood median is undefined, the comparison does not flag the pixel. We count a renderer flag only where the reference hits terrain and does not have its own flag. The result is a percentage of all image pixels, not the difference between two total flag counts.
\end{itemize}

Coverage detects gaps and added silhouettes. It cannot detect a wrong surface in front of another surface if both occupy the same pixel. Distance error helps with that case. Local incoherence detects abrupt pixel-scale distance changes, such as bands and speckle. It is a spatial measure; it does not measure temporal flicker.

The reference still differs from otherwise agreeing renderers by roughly 8--14 world units in some wide, off-axis top-down views. Halving its step did not remove that offset. The cause is not fully resolved, so we do not use absolute distance error to claim sub-unit accuracy or rank closely agreeing methods. The quality conclusions also use coverage, local incoherence, visual inspection, and agreement between independent renderers.

\subsection{Selecting the configurations}\label{sec:tuning}

We sweep each method's available quality controls on three horizon views at $400\times260$ pixels. The tuning score is missing terrain plus local incoherence. The default rule chooses the cheapest setting within one percentage point of that method's lowest score. This gives each method a useful point on its own curve; it does not give all methods equal image quality.

Table~\ref{tab:settings} lists the selected settings. Mesh receives one explicit exception. At publication resolution, the one-point rule would choose $q=0.25$, with 0.8\% combined error. We use $q=0.5$, the coarsest tested setting at the 0.3\% error floor. All reported mesh frames, memory costs, edit costs, and comparison images use that setting unless another $q$ is stated. Appendix~\ref{sec:tuning-detail} gives the sweep results and scope.

\begin{table}[tbp]\centering\small
\caption{Selected configurations. The reported comparison uses these settings throughout.}\label{tab:settings}
\begin{tabularx}{\linewidth}{@{}l Y@{}}
\toprule
Method & Setting \\
\midrule
Ray & 128 forward samples, 4 bisection steps \\
Voxel & Cell dimensions $(4,8,2)$; 100 outer and 100 inner steps \\
Sliced & 512 horizontal planes \\
Painted & No method-specific quality control \\
Scattered & Density vector $(4,4,4)$ \\
Mesh & $q=0.5$: vertical tolerance 4 world units \\
\bottomrule
\end{tabularx}
\end{table}

\subsection{Preparation, memory, and edits}\label{sec:edit-protocol}

Preparation has three CPU wall-time measures. Setup creates pipelines and uploads shared terrain data. First frame measures the first rendered frame, including work built lazily. Warmup is the total over all pre-timing frames, including the first frame. These are overlapping stages and must not be added together.

Memory accounting covers persistent data owned by each method: GPU buffers and retained CPU acceleration or fitting state. It excludes shared terrain, palette, color, depth, and shadow resources, pipelines, staging allocations, and opaque driver memory. A zero entry means no extra persistent method data, not a zero-memory renderer.

The edit removes height and upper-layer metadata in a radius-48 crater centered at $(1024,8192)$. It uses the same dirty-rectangle path as gameplay. The timing camera is at $(1024,8350)$, 40 units above the original surface, with yaw $180^\circ$ and pitch $-30^\circ$. The before/after figure uses a closer, higher view so that the crater is visible.

For each method, we construct five independent runs and take the median first post-edit frame time. We record CPU submit-and-wait time and, where available, GPU time. We compare these with ten steady frames from a fresh renderer built from the already-edited terrain.

A second check compares the updated renderer with that fresh build after 1, 2, 4, 8, and 16 frames. We call it consistent when hit/miss disagreement is at most 0.01\% of the image and the 95th-percentile distance difference is at most 0.1 world unit on shared hit pixels. This is a tolerance-based test, not a test of bitwise equality. It also reports color mean absolute error, which is not part of the pass condition.

\FloatBarrier
\section{Rendering and Editing Results}\label{sec:results}

\subsection{Views near the horizon expose sampling failures}\label{sec:pitch}

Figure~\ref{fig:quality} shows the mean quality measures at each pitch. These values come from the 780M run and were checked against the other four devices. The table in Appendix~\ref{sec:quality-detail} gives the values. Added-terrain error is small in this batch: its pitch mean is at most 0.2\% for every selected configuration.

\methodfigure{quality-pitch}{Image-quality measures by pitch on the Radeon 780M. Each value averages three scenes and is a percentage of the full image. Zero pitch looks toward the horizon. The axes use logarithmic scales; zero values are placed at 0.05\%.}{fig:quality}

Scattered has the largest coverage loss at the horizon. Its missing-terrain mean is 40.4\% of the image at $0^\circ$, compared with 0.0--0.9\% for the other methods. The three individual views range from 19.1\% to 73.2\%. Looking down reduces its coverage loss, but it still has local incoherence of 3.7\% at $-90^\circ$.

Sliced has a different failure. Its horizon coverage loss is 0.9\%, but its 2.7\% local incoherence reflects visible bands. Ray, Voxel, Painted, and the selected Mesh have much lower horizon incoherence, at 0.1--0.4\%. Painted also has a coverage loss at steeper views: 2.7\% at $-60^\circ$.

Mesh quality depends on the fit setting. In the hangar view, $q=0.0$ leaves 11.0\% of the image missing, while $q=0.5$ reduces this to about 0.5\% and restores the wall. All five devices report 0.519\% for the selected mesh in that view. The largest reported coverage spread between devices occurs in the already poor Scattered hangar result: 73.2\% on the 780M and 74.0\% on the M3.

These results show why a top-down view is an incomplete test. It hides large coverage failures in the point method and reduces the visibility of slice bands. Even at a fixed view, a method's name is not enough to describe its quality; the mesh fit setting changes whether an important wall survives.

\subsection{Mesh has the lowest mean frame time at these settings}\label{sec:timing}

Table~\ref{tab:timing} gives the mean over all twelve views. Mesh has the lowest mean on every tested device, followed by Ray. The within-session mesh and ray intervals do not overlap (Appendix~\ref{sec:timing-detail}). This ranking applies to these configurations and scenes; the methods do not all produce equal-quality images.

\begin{table}[tbp]\centering\small
\caption{Mean frame latency over twelve views, in milliseconds. Vulkan rows use GPU timestamps. The M3 row uses CPU submit-and-wait and must be read separately. Values are rounded to two decimals.}\label{tab:timing}
\setlength{\tabcolsep}{3pt}
\begin{tabular}{@{}lrrrrrr@{}}
\toprule
Device & Ray & Voxel & Sliced & Scattered & Painted & Mesh \\
\midrule
780M & 4.83 & 8.32 & 9.34 & 9.22 & 15.84 & \textbf{4.07} \\
7900 XT & 0.65 & 1.40 & 1.63 & 2.79 & 5.22 & \textbf{0.39} \\
Intel & 18.17 & 35.57 & 41.69 & 88.20 & 79.63 & \textbf{11.45} \\
RTX 5070 & 0.79 & 1.26 & 1.90 & 1.47 & 4.69 & \textbf{0.50} \\
\midrule
M3 (CPU) & 8.34 & 12.49 & 12.13 & 18.45 & 32.72 & \textbf{5.92} \\
\bottomrule
\end{tabular}
\end{table}

Voxel costs more than Ray on every device in this batch. Its empty-space skipping does not repay the hierarchy traversal and sampling costs at the selected settings. This is a result for the tested grid and views, not a general limit on voxel acceleration.

Painted becomes more expensive as the view turns down and its ground footprint grows. On Intel, its pitch mean rises from 45.6 ms at the horizon to 124.0 ms looking straight down. Scattered is sensitive to both scene and device. The hangar view takes 29.3 ms on the 780M, 21.8 ms on the 7900 XT, and 148.9 ms on Intel, but 2.5 ms on the RTX 5070. Its mean should therefore not be treated as a stable cost for all horizon views.

\subsection{Fast draws can require expensive preparation}\label{sec:prep}

Table~\ref{tab:prep} separates preparation from timed rendering. Mesh fits the terrain on the CPU in a blocking first frame of about 2.5 s on the 780M system. Voxel spreads its initial occupancy build over several frames; the warmup totals about 3.2 s. Terrain in regions that have not yet been baked may be missing during that process.

\begin{table}[tbp]\centering\small
\caption{Preparation wall time on the 780M system, in milliseconds. Each entry is the maximum over twelve scenes. Warmup includes the first frame.}\label{tab:prep}
\begin{tabular}{@{}lrrr@{}}
\toprule
Method & Setup & First frame & All warmup \\
\midrule
Ray & 9 & 60 & 101 \\
Voxel & 12 & 124 & 3163 \\
Sliced & 17 & 88 & 144 \\
Scattered & 10 & 127 & 362 \\
Painted & 18 & 119 & 200 \\
Mesh & 10 & 2464 & 2488 \\
\bottomrule
\end{tabular}
\end{table}

The initial build is a separate application cost. A steady-state timing table hides both the mesh's blocking fit and the voxel method's incomplete first frames. The best choice can change when a short level-load time matters more than a small saving on each later frame.

\subsection{Edits reveal the cost of retained state}\label{sec:edit}

All six methods accept the same crater edit without reloading the level. Figure~\ref{fig:edit} shows the three update paths: a direct texture reader, an occupancy update, and a mesh refit. Table~\ref{tab:edit} reports the first updated frame on the 780M.

\methodfigure{edit}{The same crater before the edit (top) and on the first updated frame (bottom). Ray reads the edited texture, Voxel updates occupancy, and Mesh refits chunks. These views show the shape of the edit; timing uses the camera defined in Section~\ref{sec:edit-protocol}.}{fig:edit}

\begin{table}[tbp]\centering\small
\caption{Edit latency on the 780M, in milliseconds. Steady is a fresh build of the edited terrain; edit is the first post-edit frame. CPU time includes upload, CPU update work, submission, and waiting. Consistency uses the thresholds in Section~\ref{sec:edit-protocol}.}\label{tab:edit}
\setlength{\tabcolsep}{3pt}
\begin{tabular}{@{}lrrrrl@{}}
\toprule
& \multicolumn{2}{c}{CPU time} & \multicolumn{2}{c}{GPU time} & Consistent \\
Method & Steady & Edit & Steady & Edit & after \\
\midrule
Ray & 6.51 & 8.70 & 5.59 & 7.45 & 1 frame \\
Voxel & 7.54 & 9.00 & 6.60 & 6.97 & 1 frame \\
Sliced & 7.79 & 10.97 & 6.91 & 9.60 & 1 frame \\
Scattered & 8.41 & 10.16 & 7.37 & 8.65 & 1 frame \\
Painted & 17.89 & 22.72 & 16.87 & 21.04 & 1 frame \\
Mesh & 6.07 & 15.94 & 4.77 & 6.07 & Not by 16 \\
\bottomrule
\end{tabular}
\end{table}

Mesh has the largest measured CPU increase on the 780M: 9.88 ms, compared with a GPU increase of 1.30 ms. The CPU increase ranges from 5.0 to 9.9 ms across the five devices. The delta uses the unrounded measurements, so it can differ slightly from subtracting displayed values.

The five other methods pass the fresh-build consistency test on the first updated frame on every device. Mesh does not pass by frame 16. It differs by 0.0016\% in hit/miss classification and by 0.22 world unit at the 95th percentile of distance error. Color mean absolute error is 0.78--0.80 levels on an 8-bit scale. Its hit/miss difference passes the chosen threshold, but its distance difference does not. Retaining and refining the old triangulation leaves a small, persistent difference from a fresh fit.

Table~\ref{tab:memory} shows the retained data that makes these update paths possible. The selected mesh holds 318.7 MiB of GPU geometry and 534.7 MiB of CPU triangulation. These totals include three LODs and the live fitting structures. Counting only the finest mesh would understate the cost of this editable renderer.

\begin{table}[tbp]\centering\small
\caption{Persistent method-specific data at $1280\times800$, in MiB. Shared resources, transient allocations, and driver memory are excluded.}\label{tab:memory}
\begin{tabular}{@{}lrr@{}}
\toprule
Method & GPU data & CPU data \\
\midrule
Ray & 0 & 0 \\
Voxel & 18.29 & $<0.01$ \\
Sliced & 0 & 0 \\
Scattered & 3.91 & 0 \\
Painted & $<0.01$ & 0 \\
Mesh & 318.7 & 534.7 \\
\bottomrule
\end{tabular}
\end{table}

\FloatBarrier
\section{Why the Mesh Stays Large}\label{sec:fit}

The selected mesh draws quickly but retains substantial data. We next ask what makes the fit large. Is it the rough floor, the upper layer, or both? We use two forms of evidence: a survey across worlds and a count of the samples that trigger insertion in one world.

\subsection{Quality and output size}\label{sec:fit-quality}

Table~\ref{tab:fit-quality} shows the size of the Fostral fit as its tolerance changes. The reduction ratio is $2N/T$, where $N$ is the number of source texels and $T$ is the emitted triangle count. The baseline is a grid with two triangles per source texel. This is a triangle-count ratio, not a ratio of stored bytes. Counts in the table are rounded; ratios use the unrounded counts.

\begin{table}[tbp]\centering\small
\caption{Fostral mesh size at five quality settings. The rendering comparison uses $q=0.5$; the world survey below uses $q=0.25$.}\label{tab:fit-quality}
\begin{tabular}{@{}rrrrr@{}}
\toprule
$q$ & Tolerance & Vertices & Triangles & Reduction \\
\midrule
0.0 & 16 & 1.74 M & 3.39 M & $19.8\times$ \\
0.25 & 8 & 2.31 M & 4.50 M & $14.9\times$ \\
0.5 & 4 & 4.28 M & 8.47 M & $7.9\times$ \\
0.75 & 2 & 7.14 M & 14.1 M & $4.7\times$ \\
1.0 & 1 & 11.3 M & 22.4 M & $3.0\times$ \\
\bottomrule
\end{tabular}
\end{table}

A tighter fit needs more vertices, as expected. The question is why some worlds need many more than others at the same tolerance.

\subsection{A ten-world survey}\label{sec:survey}

We fit all ten stock worlds at $q=0.25$, or an 8-unit vertical tolerance. They share the terrain encoding, height precision, texel scale, and authoring tools, but differ in content. The three worlds without an upper layer reduce by $45.3$--$182.0\times$. The seven worlds with upper layers reduce by $5.2$--$19.4\times$ (Table~\ref{tab:survey}).

\begin{table}[tbp]\centering\small
\caption{Mesh reduction at $q=0.25$ across ten worlds. Dual is the percentage of dual-level texels. Roughness is the mean absolute discrete Laplacian of the floor or topmost surface.}\label{tab:survey}
\setlength{\tabcolsep}{5pt}
\begin{tabular}{@{}lrrrr@{}}
\toprule
World & Dual (\%) & Reduction & Floor rough. & Surface rough. \\
\midrule
weexow & 0.0 & $182.0\times$ & 1.84 & 1.84 \\
ark-a-znoy & 0.0 & $154.8\times$ & 5.52 & 5.52 \\
threall & 0.0 & $45.3\times$ & 2.85 & 2.85 \\
xplo & 1.4 & $19.4\times$ & 7.54 & 8.80 \\
khox & 4.8 & $16.1\times$ & 6.94 & 7.71 \\
fostral & 10.9 & $14.9\times$ & 3.44 & 6.54 \\
necross & 17.0 & $8.0\times$ & 4.41 & 14.43 \\
glorx & 30.0 & $6.1\times$ & 6.53 & 14.66 \\
boozeena & 13.3 & $5.7\times$ & 3.44 & 25.15 \\
hmok & 38.0 & $5.2\times$ & 2.71 & 18.77 \\
\bottomrule
\end{tabular}
\end{table}

We measure roughness as the mean of the absolute four-neighbor discrete Laplacian:
\begin{equation}
R(f)=\operatorname{mean}_{i,j}
\left|4f_{i,j}-f_{i-1,j}-f_{i+1,j}-f_{i,j-1}-f_{i,j+1}\right|.
\end{equation}
The mean excludes the outer row and column. Floor roughness uses \code{low}. Surface roughness uses \code{high} in a dual-level column and \code{low} elsewhere. It measures the topmost surface, not all three fields used by the fitter.

Across these worlds, the correlation between log reduction and floor roughness is $-0.17$. It is $-0.77$ for the dual-level fraction and $-0.82$ for topmost-surface roughness. For example, hmok has a slightly smoother floor than threall by this measure (2.71 versus 2.85), yet reduces by $5.2\times$ rather than $45.3\times$.

The survey links poor reduction more strongly to layered content than to floor roughness. It does not isolate a causal effect by adding and removing slabs in the same world: other content also changes between worlds. The insertion count below gives more direct evidence of where this fitter spends vertices.

\subsection{Layer boundaries drive extra insertions}\label{sec:boundary}

For Fostral at $q=0.25$, we classify each vertex insertion by its cause (Table~\ref{tab:insertions}). The boundary between single- and dual-level terrain accounts for 23.3\% of insertions.

\begin{table}[tbp]\centering\small
\caption{Causes of vertex insertion in the Fostral fit at $q=0.25$.}\label{tab:insertions}
\begin{tabular}{@{}lr@{}}
\toprule
Cause & Share \\
\midrule
Floor field & 53.5\% \\
Slab interior & 18.8\% \\
Single-/dual-level boundary & 23.3\% \\
Chunk-border simplification & 4.4\% \\
\bottomrule
\end{tabular}
\end{table}

At a single-level sample, both \code{mid} and \code{high} equal the floor height. At a neighboring dual-level sample, they describe a ceiling and a slab top. A triangle crossing that boundary interpolates between different structural roles. The fitter reduces the resulting vertical error by inserting more points near the boundary.

Boundary insertions change relatively little across quality settings: from 338 thousand at $q=0$ to 396 thousand at $q=1$. Floor-driven insertions grow from 444 thousand to 6.1 million over the same range. For this fitter and world, the boundary creates a substantial cost even at a loose tolerance.

A triangulation constrained to the slab outline could represent that boundary more directly. It could also avoid triangles that cross between incompatible layer states. We have not implemented or measured that alternative. The result here concerns the shared, unconstrained topology used by this mesh, rather than an unavoidable cost of every layered-terrain mesh.

\FloatBarrier
\section{Choosing a Method}\label{sec:choice}

For the tested workload, Mesh and Ray offer the clearest choices. Mesh combines the lowest mean frame time with low measured horizon error at $q=0.5$. Its costs occur outside the ordinary draw: a blocking initial fit, 319 MiB of GPU geometry, 535 MiB of CPU triangulation, and local CPU refits. Its edited output also depends on the retained triangulation.

Ray is useful when memory, load time, and immediate access to edited terrain matter more than the measured frame-time gap. It keeps no derived terrain representation. Its main weakness is sampling: a fixed step budget can miss thin features along long rays. Increasing that budget trades frame time for better coverage.

Voxel has lower missing-terrain error than Ray in the selected views, but costs more frame time and needs an occupancy build. It may be useful for a different balance of empty space and surface detail; this experiment does not establish such a case. Sliced is simple to update but shows grazing bands. Painted avoids a stored mesh but becomes expensive for large ground footprints. Scattered reads edits directly, yet its single-pixel point writes leave large gaps in the horizon scenes.

An explicit mesh may also serve other parts of an engine. For example, a rigid-body system that accepts triangle colliders could reuse it. We did not measure collision quality or collider rebuild cost. Reusing a visual mesh would add a requirement to update the collider when the corresponding terrain chunk changes.

The general lesson is to choose a method using the full workload. View direction and feature size determine whether a sampling scheme is adequate. Preparation, retained memory, and edit latency determine whether the representation fits the application. A faster ordinary frame does not settle those questions.

\section{Limitations}\label{sec:limits}

The rendering comparison uses one world and one terrain format. The ten-world survey broadens the mesh-size evidence, but all ten worlds come from Vangers. The results support decisions for this data and these implementations; broader claims need other terrain sets and stronger implementations of each method family.

The quality settings are not equal-quality operating points. The sweeps use three horizon views, mostly at a smaller resolution, and Mesh has an explicit selection exception. The point scatterer has no filtered footprint. The mesh supplies a different normal input. Each choice limits how far the measured rankings can be generalized.

The CPU reference has an unresolved distance offset in some views. Coverage and local incoherence provide useful checks, but neither is a complete perceptual metric. Static images and a supplemental flythrough do not replace a quantitative temporal-stability test.

The edit experiment uses one crater shape at one location. It measures the first update and agreement with a fresh build, not sustained destruction, long edit histories, or collider maintenance. The memory totals cover known persistent payloads rather than total process or driver memory.

The hardware study covers four Vulkan devices and one Metal device. It does not cover D3D12, mobile GPUs, or multiple driver versions per device. A separate browser smoke test demonstrates execution of the Voxel route in Firefox; it is not a six-method browser benchmark. The frame intervals describe isolated frames within one session, and the M3 uses a different timing mechanism.

\section{Conclusion}\label{sec:conclusion}

The six methods expose different tradeoffs when they render the same editable layered terrain. Horizon views reveal large gaps in point scattering and bands in horizontal slicing. At the selected settings, Mesh has the lowest mean frame time on all five tested devices. Ray gives up some speed while avoiding a retained acceleration structure and reading edits on the next draw.

The mesh's main costs are its initial fit, retained state, and local refits. A survey across ten worlds links poor reduction to layered content, and the Fostral insertion count identifies slab boundaries as a major cost for the shared topology. These results point toward better boundary handling as future work.

A useful terrain comparison must report more than ordinary frame time. It should state the view and quality settings, show the resulting images, and measure the preparation, memory, and edit costs that the application will pay.

\section*{Code, Data, and Supplemental Material}\label{sec:availability}

The engine and evaluation tools are Apache-2.0 and are available from the \href{https://github.com/kvark/vange-rs}{vange-rs repository}. The measurements use tag \code{terrain-paper}, commit \code{5586e2a}. From that revision, \path{tools/compare-terrain.py} runs the default rendering and edit protocol. Separate commands run the tuning sweeps and the world survey; \path{paper/README.md} records them. The reported tables are in \path{paper/results.md}, \path{paper/tuning.md}, and \path{paper/survey.json}.

Fostral data is published by Association K-D Lab under CC BY-SA 4.0 in \href{https://github.com/KranX/Vangers/tree/f1ad7d7/data/thechain/fostral}{the Vangers source tree}, at commit \code{f1ad7d7} \citep{vangerssource}. The harness fetches that version. The nine other survey worlds require a lawfully obtained game copy. Figures derived from Fostral and the supplemental video carry the same data attribution.

The supplemental flythrough shows the six selected methods along one camera path. It is generated by \path{tools/render-paper-video.py} and is packaged as \path{anc/terrain-methods.mp4} when the video is supplied. It does not show the crater experiment.

\subsection*{Acknowledgements}

I thank Association K-D Lab for \emph{Vangers} and the published Fostral data, Yury Zhuravlev for maintaining the open-source Vangers tree, and the players and maintainers who have kept the game alive.

OpenAI Codex, Anthropic Claude, and xAI Grok assisted with drafting, editing, literature search, and the evaluation tools. The author is responsible for the manuscript and its results.

\FloatBarrier
\clearpage
\appendix
\section{Ray-March Pseudocode}\label{sec:ray-code}

Listing~\ref{lst:ray} expands the direct marcher in Section~\ref{sec:ray}. The step is computed once from the original endpoints. The endpoint test covers the final part of the clipped segment. Refinement runs only after a solid sample has been found. This still does not guarantee the first geometric intersection when a forward step crosses several solid and empty intervals.

\noindent\begin{minipage}{\linewidth}
\begin{lstlisting}[float=false,caption={Direct height-field ray marching over a clipped segment. SOLID is Equation~\ref{eq:solid}.},label={lst:ray}]
function RAY(a0, b0, steps):
    delta = (b0 - a0) / (steps + 1)
    a = a0; b = b0; hit = false
    for i = 1 to steps:
        c = a + delta
        if SOLID(c):
            b = c; hit = true; break
        a = c
    if not hit:
        hit = SOLID(b0)
    if not hit:
        return MISS
    for i = 1 to 4:
        c = (a + b) / 2
        if SOLID(c): b = c
        else: a = c
    return b
\end{lstlisting}
\end{minipage}

\section{Configuration Selection}\label{sec:tuning-detail}

The sweep score is the sum of missing terrain and local incoherence, in percentage points of the image. It excludes added-terrain and absolute distance error. Table~\ref{tab:sweeps} summarizes the small-resolution sweeps. Their times choose among settings in that sweep; they are not the five-device publication timings.

\begin{table}[tbp]\centering\small
\caption{Quality sweeps on three horizon views at $400\times260$. The rule selects the cheapest setting within one point of the method's best score.}\label{tab:sweeps}
\setlength{\tabcolsep}{4pt}
\begin{tabularx}{\linewidth}{@{}l Y l r@{}}
\toprule
Method & Values tested & Selected & Score (\%) \\
\midrule
Ray & 16--256 steps & 128 & 1.2 \\
Voxel & Two grids; 40--400 steps & $(4,8,2)$; 100 & 0.3 \\
Sliced & 32--512 planes & 512 & 6.5 \\
Painted & No quality control & Default & 0.1 \\
Scattered & Density 1--4 per axis & $(4,4,4)$ & 39.2 \\
Mesh & $q=0.0$--$1.0$ & $q=0.5$ & 1.0 \\
\bottomrule
\end{tabularx}
\end{table}

Ray's score falls from 6.9\% at 16 steps to 0.6\% at 256. The 128-step setting reaches 1.2\%, within one point of the best. Voxel reaches 0.3\% with both step budgets set to 100; 200 and 400 do not improve the reported score. Sliced improves from 16.7\% at 32 planes to 6.5\% at 512. Scattered still has a 39.2\% score at the highest tested density. Selecting its best tested point does not make it an adequate horizon renderer.

Mesh was also swept at $1280\times800$. The scores for $q=0.0$, $0.25$, and $0.5$ are 3.9\%, 0.8\%, and 0.3\%; $q=0.75$ and $1.0$ also report 0.3\%. The one-point rule would select $q=0.25$. The publication instead uses $q=0.5$ to reach the observed error floor. This is a stricter mesh quality choice, not an application of the same rule.

The selected Voxel grid uses 18.29 MiB. The production $(2,4,1)$ grid uses about 153 MiB. The comparison characterizes the coarse grid; it does not establish the best setting for every device. Changes to resolution, view distance, or scene distribution can change the selected settings.

\section{Quality Values by Pitch}\label{sec:quality-detail}

Table~\ref{tab:quality} supplies the values behind Figure~\ref{fig:quality}. Missing-terrain and incoherence percentages both use the full image as the denominator. Values are means over the three views at each pitch.

\begin{table}[htbp]\centering\small
\caption{Missing terrain / local incoherence (\%) on the 780M.}\label{tab:quality}
\setlength{\tabcolsep}{3pt}
\begin{tabular}{@{}lrrrrrr@{}}
\toprule
Pitch & Ray & Voxel & Sliced & Scattered & Painted & Mesh \\
\midrule
$0^\circ$ & 0.6 / 0.4 & 0.0 / 0.2 & 0.9 / 2.7 & 40.4 / 1.8 & 0.0 / 0.1 & 0.3 / 0.1 \\
$-30^\circ$ & 0.6 / 0.9 & 0.0 / 0.2 & 2.3 / 1.4 & 15.4 / 3.5 & 0.3 / 0.1 & 0.4 / 0.1 \\
$-60^\circ$ & 0.1 / 0.7 & 0.0 / 0.3 & 0.0 / 0.2 & 0.3 / 5.7 & 2.7 / 0.1 & 0.0 / 0.1 \\
$-90^\circ$ & 0.1 / 0.5 & 0.0 / 0.3 & 0.0 / 0.2 & 0.1 / 3.7 & 0.8 / 0.2 & 0.0 / 0.1 \\
\bottomrule
\end{tabular}
\end{table}

\section{Timing Uncertainty}\label{sec:timing-detail}

Table~\ref{tab:intervals} gives approximate within-session 95\% intervals for the twelve-view mean. For scene $j$, let $s_j$ be the sample standard deviation of its $n_j=40$ frame times. The reported half-width is $1.96\sqrt{\sum_j s_j^2/n_j}/12$. This calculation treats the samples as independent; it does not model temporal correlation. Scenes are fixed, so the intervals do not estimate uncertainty over all possible scenes. The displayed $0.00$ half-widths are values rounded below 0.005 ms, not zero uncertainty. Per-view times and device records are retained in \path{paper/results.md}.

\begin{table}[htbp]\centering\small
\caption{Mean frame latency $\pm$ within-session 95\% interval half-width, in milliseconds. M3 uses CPU timing; the other rows use GPU timing.}\label{tab:intervals}
\setlength{\tabcolsep}{2.5pt}
\begin{tabular}{@{}lrrrrrr@{}}
\toprule
Device & Ray & Voxel & Sliced & Scattered & Painted & Mesh \\
\midrule
780M & $4.83\pm.12$ & $8.32\pm.01$ & $9.34\pm.12$ & $9.22\pm.04$ & $15.84\pm.09$ & $4.07\pm.05$ \\
7900 XT & $.65\pm.01$ & $1.40\pm.00$ & $1.63\pm.01$ & $2.79\pm.00$ & $5.22\pm.00$ & $.39\pm.00$ \\
Intel & $18.17\pm.10$ & $35.57\pm.06$ & $41.69\pm.42$ & $88.20\pm.16$ & $79.63\pm.44$ & $11.45\pm.04$ \\
RTX 5070 & $.79\pm.00$ & $1.26\pm.00$ & $1.90\pm.00$ & $1.47\pm.00$ & $4.69\pm.00$ & $.50\pm.00$ \\
\midrule
M3 (CPU) & $8.34\pm.03$ & $12.49\pm.01$ & $12.13\pm.02$ & $18.45\pm.02$ & $32.72\pm.04$ & $5.92\pm.07$ \\
\bottomrule
\end{tabular}
\end{table}

\section{Implementation and Validation Notes}\label{sec:implementation}

The original Vangers software renderer projected terrain in object order using a painter's algorithm \citep{vangerssource}. The vange-rs project began in 2016 and added the GPU methods over several years. It adopted native wgpu in 2019 and moved its shaders to WGSL in 2021 \citep{malyshau2020webgpu,malyshau2021rust}. This history explains the range of methods; it is not evidence of algorithmic novelty.

A separate browser test compiled the same source to WebAssembly and ran the Voxel route in Firefox 152 on a Radeon 890M. It loaded Fostral and rendered a $1280\times714$ canvas through WebGPU. This checks that a browser path exists. The five-device measurements in this paper remain native Vulkan and Metal runs. Commands, versions, and hashes for the browser test are in \path{paper/browser-smoke.md}.

During harness development, we corrected reference-camera focal scaling, the horizontal camera basis, the far-plane convention, and the test for terrain at height zero. These faults could produce plausible images or scores, so visual agreement and camera checks were necessary before collecting the final batch. The remaining reference-distance limitation is stated in Section~\ref{sec:metrics}. Resolved harness faults are not additional rendering results.

\FloatBarrier
{\footnotesize
\bibliographystyle{jcgt}
\bibliography{references}
}

\noindent\begin{minipage}{\linewidth}
\section*{Author Contact Information}
Dzmitry Malyshau\\
Independent Researcher\\
\href{mailto:kvark@fastmail.com}{kvark@fastmail.com}\\
ORCID \href{https://orcid.org/0009-0005-6410-4276}{0009-0005-6410-4276}

\vspace{2mm}\hrule\vspace{2mm}
{\small
\noindent \copyright~2026~Malyshau (the Author).
This manuscript is provided under the Creative Commons CC BY-ND 4.0 license:
\url{https://creativecommons.org/licenses/by-nd/4.0/}.
The Author also permits reuse of images and text from the first page to promote or summarize the work in scholarly venues, with a scientific citation to the work.

\vspace{1mm}
\noindent\includegraphics[width=2.0cm]{CC-BY-ND.png}
}
\end{minipage}

\end{document}
